% ==============================================================================
% CHAPTER 5: TYPICAL APPLICATION DOMAINS
% ==============================================================================
\section{Typical Application Domains}
\label{sec:applications}

Multi-agent collaborative QA systems are rapidly being deployed across
multiple vertical sectors, capitalizing on their advantages in complex
logical processing and expert knowledge integration.
Table~\ref{tab:app_compare} summaries the distinguishing features and
representative systems in each application domain.

\begin{table*}[pos=t]
  \centering
  \caption{Comparison of typical applications of multi-agent collaborative QA systems.}
  \label{tab:app_compare}
  \small
  \setlength{\tabcolsep}{3pt}
  \begin{tabular}{p{1.8cm} p{2.4cm} p{4.0cm} p{4.0cm} p{4.0cm}}
    \toprule
    \textbf{Domain} & \textbf{Representative System} &
    \textbf{Core Challenge} & \textbf{Collaboration Pattern} &
    \textbf{Performance Highlights} \\
    \midrule
    Healthcare &
      MedAgents \citep{tang2024medagents} &
      Extremely high factual accuracy demands; long-tail rare-disease diagnosis &
      Multi-specialty physician role-play; symptom--pathology cross-verification &
      Improves zero-shot medical reasoning across nine evaluated datasets \\[6pt]
    Legal &
      ChatLaw \citep{cui2023chatlaw} &
      Retrieval from voluminous legal statutes; rigorous formal-logic reasoning &
      Plan--Retrieve--Generate pipeline; statute-matching and verdict-prediction agents collaborate &
      More rigorous responses to long case texts on LegalBench \\[6pt]
    Education &
      GenMentor \citep{li2024genmentor} &
      Heuristic guidance rather than direct answers; personalised adaptation &
      Tutor, encourager, and evaluator roles in a mixed-role network &
      Multi-path learning recommendations; interaction closer to a human private tutor \\[6pt]
    Software Engineering &
      ChatDev \citep{qian2023chatdev} &
      Cross-file logic dependencies; multi-round compilation and debug feedback &
      SOP-based waterfall; product, development, and testing roles in group-chat collaboration &
      Full closed-loop of automated writing, reviewing, and running code \\
    \bottomrule
  \end{tabular}
\end{table*}

\subsection{Medical QA}
\label{sec:medical}

The medical domain imposes extremely high standards of factual accuracy
and rigour.
MedAgents \citep{tang2024medagents} proposes a multi-disciplinary
cross-verification collaborative framework that simulates a multi-specialty
expert consultation, effectively addressing the deficiencies of single-model
reasoning on complex rare-disease cases.
MDAgents \citep{kim2024mdagents} adaptively assigns solo or group
collaboration according to estimated case complexity. Med-PaLM
\citep{singhal2023medpalm}, by contrast, is a domain-specialised medical LLM
rather than a multi-agent system; it is relevant as a possible expert backbone,
not as evidence that adaptive agent recruitment itself reduces clinical risk.

\subsection{Legal QA}
\label{sec:legal}

Legal case QA typically involves precise retrieval over voluminous case
files and statutes.
For Chinese-language legal domains, ChatLaw \citep{cui2023chatlaw}
integrates an external legal knowledge base and employs RAG for multi-hop
retrieval.
LegalBench \citep{guha2023legalbench} is a general legal-reasoning benchmark,
not evidence for a particular multi-agent topology. A directly relevant
example is LegalGraphRAG \citep{legalgraphrag2026}, which separates evidence
retrieval, source auditing, and adjudication in a graph-grounded pipeline.
This design makes evidence provenance more explicit, although benchmark
performance does not by itself establish suitability for autonomous legal use.

\subsection{Educational QA}
\label{sec:education}

In intelligent education settings, QA systems should not merely provide
correct answers; they should guide students to think independently.
The GenMentor framework \citep{li2024genmentor} introduces multiple roles
including tutor, corrector, and knowledge extender to construct
goal-oriented personalised learning paths.
Agents collaboratively adjust the depth of explanation in real time
based on student feedback, realising Socratic heuristic teaching.

\subsection{Customer Service and Office Automation}
\label{sec:enterprise}

In enterprise customer service and process automation, multi-agent QA
systems can directly integrate with back-end corporate APIs.
Empirical research by \citet{brynjolfsson2023generative_ai_work}
reports a 15\% average increase in issues resolved per hour in one firm's
deployment of a generative-AI assistant.
By delegating knowledge retrieval, sentiment management, and ticket creation
to dedicated sub-agents, the system achieves 24-hour unattended, efficient
problem resolution and fault diagnosis.

\subsection{Code Generation and Technical QA}
\label{sec:code}

Code generation is one of the most mature application areas for multi-agent
systems.
SWE-agent \citep{yang2024swe_agent} and ChatDev \citep{qian2023chatdev}
convert complex software-engineering requirements into multi-round
collaborative workflows.
On the SWE-bench benchmark \citep{jimenez2024swe_bench}, collaborative
pipelines in which an architect agent generates the framework, a programmer
agent writes the code, and a testing agent provides compilation feedback
not only answer programming questions but directly generate complete
multi-thousand-line projects.
